% !TeX root = ../../main.tex
\subsection{\texorpdfstring{From CB\textsuperscript{2} to
BARCS}{From CB2 to BARCS}}

BARCS is fundamentally a guide-level coefficient model.  Its required inputs
are a guide-by-library count matrix, the corresponding full-library totals,
and a sample design matrix; its primary output is one coefficient or contrast
test per guide.  Gene aggregation is downstream and optional.  In particular,
neither Fisher nor Stouffer combination is required to fit or interpret the
BARCS regression.  Some benchmarks below use an explicitly labeled Stouffer
summary to compare with deposited gene-level results, whereas BARCS-GC uses
the direct shared-effect Wald construction in
Equations~\eqref{eq:guideconsistency}--\eqref{eq:gcz}.

CB\textsuperscript{2} adapts the beta-binomial weighted test of
\citet{baggerly2003differential} to CRISPR pooled screens
\citep{jeong2019beta}.  Its central insight is worth preserving: sequencing
depth measures a guide proportion precisely, but it does not turn reads into
independent biological replicates.  The original method applies this insight
to a two-group contrast.  BARCS retains the same
depth-aware, extra-binomial variance principle and changes the mean model from
two group means to a design matrix.  ``Retains'' refers to the stochastic
hierarchy and the sample-level information ceiling; it does not mean that the
default logit fit reproduces the original finite-sample statistic.  The
legacy representation and asymptotic relationship are established below.

\begin{table}[ht]
\centering
\caption{What changes---and what does not---from the original
CB\textsuperscript{2} method to BARCS.}
\label{tab:cb2difference}
\small
\setlength{\tabcolsep}{5pt}
\begin{tabular}{p{0.20\textwidth}p{0.32\textwidth}p{0.38\textwidth}}
\toprule
Feature & \methodswatch{cbtwocyan} Original CB\textsuperscript{2} &
\methodswatch{barcsblue} BARCS\\
\midrule
Primary question & Difference between two groups &
Trend, adjusted association, interaction, or contrast\\
Sample design & Reference and comparison labels &
Arbitrary full-rank R model matrix\\
Mean model & Two unrelated group proportions &
$\logit(\mu_i)=\mathbf{x}_i^\top\boldsymbol{\beta}$\\
Effect & Difference or fold change between groups &
Conditional log-odds coefficient or linear contrast\\
Dispersion & Separate beta-binomial weighting within each group &
One guide-specific $\rho$ shared across the fitted design\\
Estimator & Weighted group proportions &
Feasible logistic IRLS\\
Variance principle & \multicolumn{2}{c}{Binomial sequencing plus
between-library beta-binomial heterogeneity}\\
Reference df & Group-variance df formula &
Residual $m-q$\\
Adjustment & Pairwise stratification or preprocessing &
Batch, donor, dose, time, interactions, splines, and other covariates\\
Implementation & Existing CB\textsuperscript{2} workflow &
Additive \texttt{bbreg()}, \texttt{bb\_contrast()}, and
\texttt{bb\_screen()} APIs with RcppArmadillo kernels\\
\bottomrule
\end{tabular}
\end{table}

The regression bridge was described by \citet{baggerly2004overdispersed} for
sequencing count data with multiple groups and continuous covariates.  The
method combines logistic regression with extra-binomial variation and retains
a coefficient-based $t$ test.  Here we state that construction in CRISPR-screen
notation and use the beta-binomial, library-size-dependent variance of
\citet{williams1982extrabinomial}.  An identity-link regression of raw
proportions would allow fitted values outside $(0,1)$, while an unweighted
correlation would ignore unequal library sizes.  The logit and beta-binomial
weights solve those two problems directly.
